
% ============================================================================
\clearpage
\section*{v4.8.0-r06 Source-Only: Signature of the Biquaternion Kernel and the Complex-Vector Reading}
\addcontentsline{toc}{section}{v4.8.0-r06 Source-Only: Signature of the Biquaternion Kernel and the Complex-Vector Reading}
\label{sec:v480-r06-complex-vector-reading}
% ============================================================================

\paragraph{Purpose.}
Definition~\ref{def:biquaternion} gives the kernel
\(q_n=it\cdot\mathbf 1+z_1e_1+z_2e_2+z_3e_3\) with \(z_k\in\mathbb C\) and the norm
\(|q_n|^2=-(ct)^2+|z_1|^2+|z_2|^2+|z_3|^2\), read as a Minkowski metric of signature
\((-,+,+,+)\) on which the Lorentz group acts by quaternion conjugation.  This layer computes
the signature, identifies the natural invariant of \(\mathbb C\otimes\HH\), and adopts as working
reading the one under which the kernels \((z_1,z_2,z_3)\) carry the Lorentz group.  In that reading
the symmetric triole is exactly the light-like centred configuration, and the photon sector is the
base \(S_0\cong\mathbb CP^1\) of Theorem~\ref{thm:hopf}.  Earlier layers are not rewritten.

\noindent\textbf{Source-only status.}
This v4.8.0-r06 layer is source-only.  It does not build a PDF, does not assign a new DOI,
and does not authorize a public release.  The v4.8.0 Version DOI
\href{https://doi.org/10.5281/zenodo.23365244}{10.5281/zenodo.23365244} and the Concept DOI
\href{https://doi.org/10.5281/zenodo.19210728}{10.5281/zenodo.19210728} are unchanged by it.

\subsection*{v4.8.0-r06 signature and natural norm}
\label{subsec:v480-r06-signature}

Write \(z_k=x_k+iy_k\) with \(x_k,y_k\in\mathbb R\).

\paragraph{Theoremlet v4.8.0-r06.1.}
\begin{enumerate}[label=(\roman*)]
\item The norm of Definition~\ref{def:biquaternion} is the real quadratic form
\(-t^2+\sum_k(x_k^2+y_k^2)\) of signature \((1,6)\), not \((1,3)\).
\item Conjugation \(q\mapsto uqu^{-1}\) by \(u\in\HH^\times\) fixes the scalar part and acts by
\(SO(3)\) on \((x_k)\) and on \((y_k)\) separately; it contains no boost.
\item The natural invariant of \(\mathbb C\otimes\HH\) is the complex-bilinear norm
\[
N(q):=q\bar q=-t^2+z_1^2+z_2^2+z_3^2,\qquad
\operatorname{Re}N=-t^2+|x|^2-|y|^2,\quad \operatorname{Im}N=2\,x\cdot y,
\]
where \(\bar q\) is quaternion conjugation without complex conjugation; \(N(pq)=N(p)N(q)\).
\end{enumerate}

\paragraph{Proof sketch.}
(i) The Hessian of the form has eigenvalues \(-1\) once and \(+1\) six times.  (ii) For real
\(u\), \(u e_k u^{-1}\) is a rotation of the imaginary units and \(u\mathbf 1u^{-1}=\mathbf 1\).
(iii) Direct expansion of \(q\bar q\); the vector part vanishes because the coefficients commute.
Multiplicativity is inherited from \(\HH\).  \(\square\)

Hence the Minkowski reading of Definition~\ref{def:biquaternion} and the Lorentz statement of the
subsequent remark do not hold as written, and Theorem~\ref{thm:lorentz} establishes invariance
under the involution \(J\), not under the Lorentz group.

\subsection*{v4.8.0-r06 two admissible readings}
\label{subsec:v480-r06-readings}

\paragraph{Reading A (four-vector).}
On the real slice \(\{q:q^\dagger=-q\}\), where \(q^\dagger\) combines quaternion and complex
conjugation, the \(z_k\) are real and \(N(q)=-t^2+x_1^2+x_2^2+x_3^2\) is exactly Minkowski
\((1,3)\).  The map \(q\mapsto uqu^\dagger\) with \(N(u)=1\) preserves the slice and \(N\), and
changes the time component: these are genuine boosts.  The kernels then carry no phase, so the
\(120^\circ\) triole has no place.

\paragraph{Reading B (complex three-vector).}
The kernels form \(z=(z_1,z_2,z_3)\in\mathbb C^3\), and \(q\mapsto uqu^{-1}\) with
\(u\in\mathbb C\otimes\HH\), \(N(u)=1\), fixes the scalar part and acts on \(z\) by
\(SO(3,\mathbb C)\).  Since \(SO(3,\mathbb C)\cong PSL(2,\mathbb C)\cong SO^+(1,3)\), this is the
proper orthochronous Lorentz group.  The invariant is \(s(z):=z_1^2+z_2^2+z_3^2\).  The conjugation
formula of the remark after Definition~\ref{def:biquaternion} becomes correct once \(u\) is taken
complex.  Formally, \(z\) is the Riemann--Silberstein vector \(E+iB\), and
\(s=E^2-B^2+2i\,E\cdot B\) collects the two field invariants.

\paragraph{Decision.}
Reading~B is adopted as the working reading of the kernel (author's decision, 2026-10-19).
Reading~A remains available as the four-vector structure of \(\mathbb C\otimes\HH\); the two
are not exclusive.  In Reading~B the scalar \(t\) is an invariant parameter, not the time
coordinate of a four-vector; its role is open.

\subsection*{v4.8.0-r06 the triole is the light-like centred configuration}
\label{subsec:v480-r06-triole-lemma}

\paragraph{Theoremlet v4.8.0-r06.2.}
For \((z_1,z_2,z_3)\in\mathbb C^3\),
\[
z_1+z_2+z_3=0\ \text{ and }\ z_1^2+z_2^2+z_3^2=0
\quad\Longleftrightarrow\quad
(z_1,z_2,z_3)=c\,(1,\omega,\omega^2)\ \text{up to order},
\]
with \(\omega=e^{2\pi i/3}\) and \(c\in\mathbb C\).

\paragraph{Proof sketch.}
With elementary symmetric functions \(e_1,e_2,e_3\) and \(p_2=\sum z_k^2\) one has
\(p_2=e_1^2-2e_2\).  Hence \(e_1=p_2=0\) gives \(e_2=0\), and the \(z_k\) are the roots of
\(\lambda^3-e_3\), i.e.\ the three cube roots of \(e_3\).  Conversely \(1+\omega+\omega^2=0\) and
\(1+\omega^2+\omega^4=0\).  \(\square\)

The triole condition of Theorem~\ref{thm:unity}, read with equal moduli and \(120^\circ\) phases,
is therefore not only a phase condition: in Reading~B it is equivalent to the Lorentz-invariant
null condition \(s=0\) together with centring \(\sum z_k=0\).

\subsection*{v4.8.0-r06 invariant sectors}
\label{subsec:v480-r06-sectors}

\paragraph{Theoremlet v4.8.0-r06.3.}
The orbits of \(SO(3,\mathbb C)\) on \(\mathbb C^3\) are the level sets \(\{s=c\}\) for \(c\neq0\),
the punctured null cone \(\{s=0\}\setminus\{0\}\), and \(\{0\}\).  The duality rotation
\(z\mapsto iz\) maps \(s\mapsto-s\) and preserves the null cone.

\paragraph{Proof sketch.}
For \(s\neq0\) put \(e_1=z/\sqrt s\), complete it to a complex-orthonormal basis with determinant
\(1\), and rotate \(z\) to \((\sqrt s,0,0)\); this was checked constructively.  Null vectors are
treated through Theoremlet r06.4 below.  \(s(iz)=-s(z)\) is immediate.  \(\square\)

\paragraph{Sector reading (conjecture).}
The classification of \S\ref{sec:unitary} by the sign of \(\operatorname{Re}z\) is not
Lorentz-invariant in Reading~B.  A Lorentz-invariant candidate is: \(s=0\) null (photon-type);
\(s\) real positive, electric-dominated; \(s\) real negative, magnetic-dominated.  The duality
\(z\mapsto iz\) exchanges the two massive types and fixes the null sector, matching the reading of
the tachyon as \(J\)-dual of matter and of the photon as fix-sector.  The assignment
electric \(\leftrightarrow\) matter, magnetic \(\leftrightarrow\) tachyon is a model decision and is
not made here; the meaning of non-real \(s\) is open.

\subsection*{v4.8.0-r06 the null cone and the photon base}
\label{subsec:v480-r06-null-cone}

For \(\psi=(\psi_1,\psi_2)\in\mathbb C^2\) put
\(z(\psi):=(\psi_1^2-\psi_2^2,\ i(\psi_1^2+\psi_2^2),\ -2\psi_1\psi_2)\).

\paragraph{Theoremlet v4.8.0-r06.4.}
\(z(\psi)\) lies on the null cone, every non-zero null vector has exactly the two preimages
\(\pm\psi\), and the projectivized null cone \(\{[z]\in\mathbb CP^2:s(z)=0\}\) is a smooth conic
isomorphic to \(\mathbb CP^1\cong S^2\).  Under the identification of Theorem~\ref{thm:hopf} this is
the photon base \(S_0\).  The map is equivariant: \(SL(2,\mathbb C)\) on \(\psi\) induces
\(SO(3,\mathbb C)\) on \(z\).

\paragraph{Proof sketch.}
\(s(z(\psi))=(a-b)^2-(a+b)^2+4ab=0\) with \(a=\psi_1^2\), \(b=\psi_2^2\).  For a null \(z\), set
\(\psi_1^2=\tfrac12(z_1-iz_2)\), \(\psi_2^2=\tfrac12(-z_1-iz_2)\); the null condition gives
\(z_3^2=4\psi_1^2\psi_2^2\), which fixes the relative sign, so the preimage is \(\pm\psi\).  The
conic is the Veronese image of \(\mathbb CP^1\).  Equivariance and invariance of the null cone
were checked numerically for random unit biquaternions.  \(\square\)

\paragraph{Consequences.}
The symmetric triole \((1,\omega,\omega^2)\) is the image of the spinor
\(\psi\approx(0.9744+0.1283i,\ 0.3097+0.4036i)\).  For the null axis of
Theorem~\ref{thm:fractal}, the natural Lorentz-invariant definition in Reading~B is \(s=0\), not
\(|q_n|^2=0\); its directions then form exactly \(S_0\).  The transport statement
\(\KC(L_n)=L_{n-1}\) of that theorem is not addressed and remains as stated there.

\paragraph{Information-level reading (conjecture).}
A normalized \(\psi\) is a qubit; its Bloch vector is the Hopf image \(S^3\to S^2\), and
\(k(\psi)=(|\psi|^2,\psi^\dagger\sigma\psi)\) is a light-like four-vector.  Every future time-like
vector is a sum \(k(\psi_a)+k(\psi_b)\) (checked constructively).  This supports the author's
conjecture that space-time is a projection of an information level, with single spinors as
light-like and spinor pairs as massive readouts, in line with
\(S_0=\mathbb CP^1\) and \(S_1=\HH P^1\) of Theorem~\ref{thm:hopf}.  No derivation of space-time
points is claimed.

\subsection*{v4.8.0-r06 status of the affected statements}
\label{subsec:v480-r06-status}

\begin{center}
\small
\begin{tabular}{@{}>{\raggedright\arraybackslash}p{0.34\textwidth}>{\raggedright\arraybackslash}p{0.58\textwidth}@{}}
\toprule
statement & status after r06\\
\midrule
Def.~\ref{def:biquaternion}, Minkowski reading & signature \((1,6)\); replaced by \(N(q)\) and Reading~B\\
Lorentz action by quaternion conjugation & correct only with complex \(u\), \(N(u)=1\)\\
Thm.~\ref{thm:lorentz} & shows \(J\)-invariance; Lorentz invariance supplied by r06 (Reading~B)\\
Sector classification, \S\ref{sec:unitary} & not invariant; invariant candidate via \(s\) (conjecture)\\
Triole condition, Thm.~\ref{thm:unity} & equivalent to centred null configuration (r06.2)\\
Null axis, Thm.~\ref{thm:fractal} & invariant definition \(s=0\); directions \(=S_0\) (r06.4)\\
Prop.~\ref{prop:orthogonality-third-kernel} & unaffected\\
\bottomrule
\end{tabular}
\end{center}

\subsection*{v4.8.0-r06 ledger and guard update}
\label{subsec:v480-r06-ledger-guard}

\[
R_{\rm kern}^{v480\text{-}r06}=(0,0,0,0,1,1)
\]
for (signature computed; Lorentz action supplied in Reading~B; triole lemma; null cone \(=S_0\);
sector assignment not chosen; role of \(t\) open).

The layer spends no guard:
\[
R_{\rm glob}^{(v480\text{-}r06,+)}=(0,1,1,1,1),
\qquad
\Gamma_{\rm guard}^{v480\text{-}r06}=(1,1,1,1,1),
\qquad
\beta_4=1/256\ \text{unspent}.
\]
It does not imply exact finite \(C5\) zero closure, exact all-readout terminality,
\(C6/T2\), a global photon-clock/RH theorem, or global all-readout closure.

